% conclusion.tex -- 3 to 5 sentences, same numbers as the paper, limits in one sentence, future work in one sentence, no new claim.
% Placeholders only: no result exists.

\todo{RESULT, sentence 1: the question and the measurement, stated as in the contribution paragraph, scoped to one pretrained model and one data set.}

\todo{RESULT, sentence 2: the headroom outcome and the primary contrast per corruption, with the number, interval and preregistered wording.}

\todo{RESULT, sentence 3: the calibration result (realised false-alarm proportion and the verdict) and the mechanism and detection controls.}

\todo{RESULT, sentence 4: limits in one sentence (one model, one data set, narrow guarantee, no re-infill or caption corruption in the main run).}

\todo{AUTHOR, sentence 5: future work in one sentence (a second backbone; the re-infill and caption-corruption tiers; a finetuned variant of the observed-position score).}
